% ============================================================
%  News Credibility Monitor — End-Semester Report
%  Overleaf-ready  •  pdfLaTeX / XeLaTeX / LuaLaTeX
% ============================================================
\documentclass[12pt, a4paper]{article}

% ── Packages ──────────────────────────────────────────────
\usepackage[utf8]{inputenc}
\usepackage[T1]{fontenc}
\usepackage{lmodern}
\usepackage[margin=1in]{geometry}
\usepackage{graphicx}
\usepackage{booktabs}
\usepackage{array}
\usepackage{longtable}
\usepackage{enumitem}
\usepackage{amsmath}
\usepackage{hyperref}
\usepackage{xcolor}
\usepackage{listings}
\usepackage{caption}
\usepackage{subcaption}
\usepackage{float}
\usepackage{setspace}
\usepackage{titlesec}
\usepackage{fancyhdr}

% ── Colours & Links ───────────────────────────────────────
\definecolor{linkblue}{HTML}{2563EB}
\hypersetup{
    colorlinks=true,
    linkcolor=linkblue,
    citecolor=linkblue,
    urlcolor=linkblue,
}

% ── Code listing style ────────────────────────────────────
\definecolor{codebg}{HTML}{F5F5F5}
\definecolor{codeframe}{HTML}{D4D4D4}
\lstset{
    backgroundcolor=\color{codebg},
    frame=single,
    rulecolor=\color{codeframe},
    basicstyle=\ttfamily\small,
    breaklines=true,
    captionpos=b,
}

% ── Header / Footer ──────────────────────────────────────
\pagestyle{fancy}
\fancyhf{}
\fancyhead[L]{\small News Credibility Monitor}
\fancyhead[R]{\small\thepage}
\renewcommand{\headrulewidth}{0.4pt}

% ── Section formatting ───────────────────────────────────
\titleformat{\section}{\Large\bfseries}{\thesection}{1em}{}
\titleformat{\subsection}{\large\bfseries}{\thesubsection}{1em}{}
\titleformat{\subsubsection}{\normalsize\bfseries}{\thesubsubsection}{1em}{}

% ── Spacing ──────────────────────────────────────────────
\onehalfspacing

% ════════════════════════════════════════════════════════════
\begin{document}

% ── Title Page ───────────────────────────────────────────
\begin{titlepage}
    \centering
    \vspace*{2cm}

    {\Huge\bfseries News Credibility Monitor\\[0.3cm]}
    {\Large An AI-Powered Multi-Agent System for\\Automated News Verification}

    \vspace{1.5cm}

    {\large
    % ── Replace with your details ──
    \textbf{Course:} Generative AI Capstone Project\\[0.3cm]
    \textbf{Institution:} Newton School of Technology\\[0.3cm]
    \textbf{Date:} April 2026
    }

    \vspace{2cm}

    \begin{abstract}
    \noindent
    This report presents the design, implementation, and deployment of the
    \textbf{News Credibility Monitor}---an AI-powered system that assesses the
    credibility of news articles through a multi-layered pipeline combining
    classical machine learning, Retrieval-Augmented Generation (RAG), and
    multi-agent reasoning. The system employs a TF-IDF vectorizer paired with
    Logistic Regression for an initial classification signal, an in-memory
    vector store loaded from precomputed embeddings for evidence retrieval,
    and three specialised LLM-based reasoning agents whose outputs are
    synthesised by a final judge agent. The complete pipeline is orchestrated
    as a LangGraph StateGraph with conditional branching based on ML
    confidence. The frontend is a React application deployed on Vercel, and
    the backend is a FastAPI service deployed on Render. Evaluation on the
    ISOT Fake News Dataset yields 98.88\% classification accuracy on the ML
    layer, while the agentic pipeline provides interpretable, multi-perspective
    reasoning for each verdict.
    \end{abstract}

    \vfill
\end{titlepage}

% ── Table of Contents ────────────────────────────────────
\newpage
\tableofcontents
\newpage

% ════════════════════════════════════════════════════════════
\section{Introduction}
% ════════════════════════════════════════════════════════════

The proliferation of misinformation in digital media has created an urgent
need for automated tools that can assess the credibility of news content at
scale. Traditional approaches---either purely statistical classifiers or
human fact-checking---each have significant limitations: classifiers produce
opaque binary labels with no reasoning, while human review cannot keep pace
with the volume of online content.

This project addresses both shortcomings by designing a \textbf{hybrid
intelligence pipeline} that combines three complementary techniques:

\begin{enumerate}[leftmargin=2em]
    \item \textbf{Classical Machine Learning} --- A TF-IDF + Logistic
          Regression classifier provides a fast, high-accuracy initial
          signal (prediction and confidence score).
    \item \textbf{Retrieval-Augmented Generation (RAG)} --- Semantically
          similar articles from a verified dataset are retrieved to ground
          the analysis in evidence, reducing hallucination risk.
    \item \textbf{Multi-Agent Reasoning} --- Three LLM-based agents with
          distinct analytical perspectives (conservative, sceptical, neutral)
          independently analyse the article, and a final judge agent
          synthesises their outputs into a transparent, explainable verdict.
\end{enumerate}

The pipeline is orchestrated using \textbf{LangGraph}, a state-machine
framework for composable AI workflows, with a conditional branch that
bypasses the RAG retrieval step when the ML model's confidence exceeds a
configurable threshold (85\%). This design balances latency and
thoroughness: high-confidence predictions are fast-tracked, while
ambiguous cases receive additional evidence before reasoning begins.

The system is fully deployed with a React frontend on Vercel and a FastAPI
backend on Render, making it accessible as a web application for
demonstration and evaluation.


% ════════════════════════════════════════════════════════════
\section{System Overview}
% ════════════════════════════════════════════════════════════

The News Credibility Monitor is structured as a two-tier web application
with a clear separation between the presentation layer and the analysis
engine.

\subsection{High-Level Architecture}

\begin{itemize}[leftmargin=2em]
    \item \textbf{Frontend} --- A single-page React application (Vite,
          Tailwind CSS~4, React~19) deployed on Vercel. It collects the
          user's article text, displays an animated processing pipeline
          during analysis, and renders the structured verdict with agent
          reasoning, evidence summaries, and risk indicators.
    \item \textbf{Backend} --- A FastAPI service deployed on Render's free
          tier. It exposes two endpoints: a health-check (\texttt{GET /})
          and the analysis endpoint (\texttt{POST /analyze}). All ML
          inference, RAG retrieval, and LLM orchestration occur server-side.
    \item \textbf{Keep-Alive} --- A GitHub Actions cron job pings the
          Render service every 14~minutes to prevent the free-tier instance
          from entering a sleep state.
\end{itemize}

\subsection{Data Flow Summary}

The end-to-end data flow for a single analysis request proceeds as follows:

\begin{enumerate}[leftmargin=2em]
    \item User pastes an article (\(\geq\)50 words) into the web interface.
    \item The frontend sends a \texttt{POST /analyze} request to the backend.
    \item The backend preprocesses the text and runs the LangGraph pipeline.
    \item The structured result (verdict, confidence, agent reasoning,
          evidence, risk factors) is returned as JSON.
    \item The frontend normalises the response and renders results with
          staggered reveal animations.
\end{enumerate}

\begin{figure}[H]
    \centering
    \includegraphics[width=0.95\textwidth]{system_design.png}
    \caption{High-level system architecture showing the data flow from user
             input through the ML, RAG, and multi-agent layers to the final
             verdict.}
    \label{fig:architecture}
\end{figure}


% ════════════════════════════════════════════════════════════
\section{Methodology}
% ════════════════════════════════════════════════════════════

\subsection{Text Preprocessing}

All article text passes through a shared cleaning pipeline before entering
either the ML or RAG pathways:

\begin{enumerate}[leftmargin=2em]
    \item \textbf{Dateline stripping} --- Publisher attribution patterns
          (e.g., \texttt{WASHINGTON (Reuters) ---}) are removed from the
          first 50~characters only, preventing accidental content loss on
          user-supplied text.
    \item \textbf{Lowercasing} and removal of non-alphabetic characters.
    \item \textbf{Stopword removal} using NLTK's English stopword list.
\end{enumerate}

This conservative cleaning strategy preserves the semantic content of
user-provided articles while normalising the surface form for both the
TF-IDF vectorizer and the RAG retrieval query.


% ─── 3.2  ML Model ───────────────────────────────────────
\subsection{Machine Learning Classification Layer}

The ML layer provides the system's initial, low-latency credibility signal.

\subsubsection{Feature Extraction}

Articles are transformed into numerical feature vectors using a
\textbf{TF-IDF (Term Frequency--Inverse Document Frequency)} vectorizer
with the following configuration:

\begin{itemize}[leftmargin=2em]
    \item \texttt{max\_features = 10{,}000}
    \item \texttt{ngram\_range = (1, 2)} --- captures both unigrams and
          bigrams
\end{itemize}

Stopword removal is delegated to the upstream text cleaner to avoid
double-filtering.

\subsubsection{Classifier}

The classifier is a \textbf{Logistic Regression} model with the following
hyperparameters:

\begin{itemize}[leftmargin=2em]
    \item \texttt{max\_iter = 1000}
    \item \texttt{solver = lbfgs}
    \item \texttt{C = 1.0} (default regularisation)
    \item \texttt{class\_weight = balanced} --- compensates for any
          Real/Fake class imbalance in the training data
\end{itemize}

\subsubsection{Training Data}

The model is trained on the \textbf{ISOT Fake News Dataset}, which
comprises approximately 44,000 articles split between verified real news
(sourced from Reuters) and fake news (sourced from unreliable websites
flagged by Politifact). The dataset is split 80/20 for training and
evaluation (\texttt{random\_state=42}).

\subsubsection{Output}

The ML node produces two values that propagate through the rest of the
pipeline:
\begin{itemize}[leftmargin=2em]
    \item \textbf{Prediction}: \texttt{REAL} or \texttt{FAKE}
    \item \textbf{Confidence}: The classifier's predicted probability
          (0--100\%), used for conditional routing
\end{itemize}


% ─── 3.3  RAG System ─────────────────────────────────────
\subsection{Retrieval-Augmented Generation (RAG) Layer}

The RAG layer enriches the analysis with evidence from a verified reference
corpus, reducing the LLM agents' tendency to hallucinate.

\subsubsection{Embedding Strategy}

Embeddings are \textbf{precomputed offline} using the
\texttt{all-MiniLM-L6-v2} model from SentenceTransformers. A stratified
sample of 5{,}000 articles from the ISOT dataset is embedded and serialised
as a pickle file (\texttt{embeddings.pkl}). This design decision has two
key advantages:

\begin{enumerate}[leftmargin=2em]
    \item \textbf{No PyTorch at runtime} --- The production server does not
          need the 2\,GB+ PyTorch/SentenceTransformers stack, keeping memory
          usage within Render's 512\,MiB free-tier limit.
    \item \textbf{No persistent disk} --- Embeddings are loaded into an
          \textbf{ephemeral ChromaDB client} (in-memory) at startup,
          eliminating dependency on persistent storage.
\end{enumerate}

\subsubsection{Query-Time Retrieval}

At query time, the article text is cleaned and passed to ChromaDB's
built-in \textbf{DefaultEmbeddingFunction} (ONNX-based, lightweight), which
generates the query embedding without loading PyTorch. The top-$k$
($k{=}5$) most similar articles are retrieved along with their metadata
(label, subject, source) and distance scores.

\subsubsection{Conditional RAG}

The RAG retrieval step is \textbf{conditionally executed} based on the ML
model's confidence score:

\[
\text{next\_node} =
\begin{cases}
    \texttt{rag\_node}      & \text{if } \mathrm{confidence} < 85\% \\
    \texttt{agent\_a\_node} & \text{if } \mathrm{confidence} \geq 85\%
\end{cases}
\]

When the ML model is highly confident, the system bypasses RAG to reduce
latency. When confidence is low, retrieved evidence provides critical
context for the reasoning agents.


% ─── 3.4  Multi-Agent System ─────────────────────────────
\subsection{Multi-Agent Reasoning System}

The multi-agent layer is the core differentiator of this system. Rather
than relying on a single LLM call, three agents with distinct analytical
personas independently evaluate the article, and a judge agent aggregates
their outputs.

\subsubsection{Agent Topology}

The agents are orchestrated as a \textbf{LangGraph StateGraph} with the
following topology:

\begin{lstlisting}[caption={LangGraph pipeline topology},label={lst:topology}]
preprocess_node -> ml_node -> [conditional]
    |-- confidence >= 85% --> agent_a_node
    |-- confidence <  85% --> rag_node --> agent_a_node
                                              |
                                         agent_b_node
                                              |
                                         agent_c_node
                                              |
                                          judge_node
                                              |
                                         output_node --> END
\end{lstlisting}

Agents execute \textbf{sequentially} (A~$\to$~B~$\to$~C), not in parallel.
Each agent receives the full article text, the ML prediction with
confidence, and any retrieved RAG evidence.

\subsubsection{Agent Personas}

Each agent is implemented as a structured LLM prompt with a distinct role
description:

\begin{table}[H]
\centering
\caption{Multi-agent persona configuration}
\label{tab:agents}
\begin{tabular}{@{}llll@{}}
\toprule
\textbf{Agent} & \textbf{Role} & \textbf{Persona} & \textbf{LLM Provider} \\
\midrule
Agent A & Conservative  & Trusts ML signal and evidence heavily
        & Groq \\
Agent B & Sceptical     & Challenges evidence quality, highlights uncertainty
        & Groq \\
Agent C & Neutral       & Balanced, objective reasoning
        & Groq \\
\bottomrule
\end{tabular}
\end{table}

All three agents use \textbf{Llama~3.3-70B-Versatile} via the Groq API
(temperature~$= 0.3$, max tokens~$= 1{,}500$). Each produces a structured
output:

\begin{itemize}[leftmargin=2em]
    \item \textbf{Verdict}: \texttt{REAL} or \texttt{FAKE}
    \item \textbf{Confidence}: 0--100
    \item \textbf{Reasoning}: 3--5 sentence justification
\end{itemize}


% ─── 3.5  Final Judge ────────────────────────────────────
\subsection{Final Judge Agent}

The judge agent receives all three agent responses, the ML signal, and the
agreement statistics (computed before the prompt is assembled). It uses a
different LLM---\textbf{Llama~3.3-70B-Instruct} via the NVIDIA AI
Endpoints API (temperature~$= 0.2$, top-$p = 0.7$, max tokens~$=
1{,}024$)---to provide an independent synthesis.

Using a separate LLM provider for the judge reduces the risk of systematic
bias from a single model. If the NVIDIA endpoint fails, the system falls
back to the Groq API.

The judge produces:
\begin{itemize}[leftmargin=2em]
    \item \textbf{Final Verdict}: \texttt{REAL} or \texttt{FAKE}
    \item \textbf{Final Confidence}: 0--100
    \item \textbf{Consensus Summary}: Synthesis of multi-agent reasoning
    \item \textbf{Dominant Agent}: Which agent was most influential
    \item \textbf{Conflict Resolution}: How disagreements were handled
\end{itemize}


% ─── 3.6  Output & Risk ─────────────────────────────────
\subsection{Output Assembly and Risk Detection}

The final \texttt{output\_node} assembles the structured response and
computes several \textbf{risk factors}:

\begin{itemize}[leftmargin=2em]
    \item ML confidence below 70\%
    \item Mixed labels in retrieved RAG evidence
    \item Weak agent consensus (Low agreement level)
    \item Insufficient evidence (fewer than 2 retrieved documents)
    \item Repetitive retrieved content (Jaccard similarity $> 0.4$ between
          top-2 documents)
\end{itemize}

These risk indicators are surfaced in the frontend to help users calibrate
their trust in the verdict.


% ─── 3.7  Agent Workflow ─────────────────────────────────
\subsection{Agent Workflow}

The end-to-end agent workflow follows a structured, multi-stage pipeline
orchestrated by LangGraph:

\begin{enumerate}[leftmargin=2em]
    \item \textbf{ML Classification} --- The cleaned article is vectorised
          and classified by the Logistic Regression model, producing a
          prediction (\texttt{REAL}/\texttt{FAKE}) and a confidence score.
    \item \textbf{Conditional RAG} --- If ML confidence falls below the
          85\% threshold, the system retrieves the top-5 semantically
          similar articles from the in-memory ChromaDB store. If confidence
          is high, this step is skipped to reduce latency.
    \item \textbf{Multi-Agent Reasoning} --- Three agents execute
          sequentially (A~$\to$~B~$\to$~C). Each agent receives the full
          article text, the ML prediction with confidence, and any
          retrieved evidence documents. Despite sequential execution, each
          agent reasons independently---its prompt contains only the shared
          evidence, not the outputs of preceding agents---ensuring that
          each perspective is uncontaminated by the others.
    \item \textbf{Judge Aggregation} --- The judge agent receives all
          three agent responses together with the ML signal and
          pre-computed agreement statistics, and synthesises a final
          verdict.
    \item \textbf{Output Assembly} --- The output node compiles the
          structured response, computes risk factors, and returns the
          result to the API caller.
\end{enumerate}

This staged design ensures that each component operates on well-defined
inputs and that failures at any node are isolated and recoverable (every
node returns fallback values on exception).


% ─── 3.8  Prompting Strategy ─────────────────────────────
\subsection{Prompting Strategy}

Prompt engineering is central to the quality and consistency of agent
outputs. The system uses a shared prompt template with role-specific
preambles:

\begin{itemize}[leftmargin=2em]
    \item \textbf{Role injection} --- Each agent prompt begins with a
          persona description that shapes reasoning style:
          \begin{itemize}
              \item Agent~A is instructed to trust the ML signal and
                    retrieved evidence heavily.
              \item Agent~B is instructed to challenge evidence quality and
                    highlight uncertainty.
              \item Agent~C is instructed to weigh all signals objectively
                    without strong bias.
          \end{itemize}
    \item \textbf{Evidence grounding} --- The prompt includes the article
          text (truncated to 1{,}500 characters), the ML prediction string,
          and formatted RAG evidence (label, subject, source, and a
          300-character snippet per document).
    \item \textbf{Structured output enforcement} --- Agents are required to
          respond in exactly three labelled fields:
          \texttt{Verdict}, \texttt{Confidence}, and \texttt{Reasoning}.
          This deterministic format enables reliable regex-based parsing in
          the output node.
    \item \textbf{Anti-hallucination directive} --- The system message
          explicitly states: \textit{``Respond ONLY based on the evidence
          provided. Never invent or hallucinate facts.''} This constraint
          anchors reasoning to the supplied ML and RAG signals.
\end{itemize}

The judge prompt follows a similar structure but requests additional fields
(\texttt{Consensus Summary}, \texttt{Dominant Agent},
\texttt{Conflict Resolution}) and receives the pre-computed agreement
level and verdict distribution to aid synthesis.


% ─── 3.9  State Management ───────────────────────────────
\subsection{State Management}

Pipeline state is managed through a \textbf{LangGraph \texttt{AgentState}
TypedDict} that is passed between all nodes. The state dictionary contains
the following fields:

\begin{table}[H]
\centering
\caption{AgentState fields and their lifecycle}
\label{tab:state}
\begin{tabular}{@{}llp{5.5cm}@{}}
\toprule
\textbf{Field} & \textbf{Type} & \textbf{Populated by} \\
\midrule
\texttt{article\_text}      & \texttt{str}           & Caller (initial input) \\
\texttt{cleaned\_text}      & \texttt{str}           & \texttt{preprocess\_node} \\
\texttt{ml\_prediction}     & \texttt{str}           & \texttt{ml\_node} \\
\texttt{ml\_confidence}     & \texttt{float}         & \texttt{ml\_node} \\
\texttt{retrieved\_docs}    & \texttt{List[dict]}    & \texttt{rag\_node} (conditional) \\
\texttt{agent\_a\_response} & \texttt{str}           & \texttt{agent\_a\_node} \\
\texttt{agent\_b\_response} & \texttt{str}           & \texttt{agent\_b\_node} \\
\texttt{agent\_c\_response} & \texttt{str}           & \texttt{agent\_c\_node} \\
\texttt{judge\_response}    & \texttt{str}           & \texttt{judge\_node} \\
\texttt{final\_report}      & \texttt{dict}          & \texttt{output\_node} \\
\texttt{error}              & \texttt{Optional[str]} & Any node (on failure) \\
\bottomrule
\end{tabular}
\end{table}

Each node receives the full state dictionary and returns a \emph{partial}
update containing only the fields it populates. LangGraph merges the
returned partial into the running state automatically, ensuring that
downstream nodes always see a consistent, progressively enriched view of
the data. This merge-based design eliminates manual state-passing logic
and guarantees that no node can accidentally overwrite fields it does not
own.


% ════════════════════════════════════════════════════════════
\section{System Architecture}
% ════════════════════════════════════════════════════════════

\subsection{Technology Stack}

\begin{table}[H]
\centering
\caption{Technology stack summary}
\label{tab:stack}
\begin{tabular}{@{}lll@{}}
\toprule
\textbf{Layer}    & \textbf{Technology}                  & \textbf{Deployment} \\
\midrule
Frontend          & React 19, Vite 8, Tailwind CSS 4     & Vercel \\
Backend API       & FastAPI, Uvicorn                      & Render (free tier) \\
ML Model          & scikit-learn (Logistic Regression)    & Loaded via joblib \\
Vector Store      & ChromaDB (Ephemeral Client)           & In-memory at runtime \\
Embeddings        & all-MiniLM-L6-v2 (precomputed)       & Pickle file \\
Agent Orchestration & LangGraph StateGraph                & 8-node pipeline \\
LLM (Agents)      & Llama 3.3-70B-Versatile              & Groq API \\
LLM (Judge)       & Llama 3.3-70B-Instruct               & NVIDIA AI Endpoints \\
Keep-Alive        & GitHub Actions (cron)                 & Every 14 minutes \\
\bottomrule
\end{tabular}
\end{table}

% ════════════════════════════════════════════════════════════
\section{Results and Analysis}
% ════════════════════════════════════════════════════════════

\subsection{ML Model Performance}

The Logistic Regression classifier was evaluated on a held-out 20\% test
split of the ISOT dataset.

\begin{table}[H]
\centering
\caption{ML model evaluation metrics on the ISOT test set}
\label{tab:ml-metrics}
\begin{tabular}{@{}lc@{}}
\toprule
\textbf{Metric}              & \textbf{Score} \\
\midrule
Accuracy                     & 98.88\% \\
Precision (weighted)         & 98.88\% \\
Recall (weighted)            & 98.88\% \\
F1 Score (weighted)          & 98.88\% \\
\bottomrule
\end{tabular}
\end{table}

\begin{table}[H]
\centering
\caption{Per-class performance}
\label{tab:per-class}
\begin{tabular}{@{}lccc@{}}
\toprule
\textbf{Class} & \textbf{Precision} & \textbf{Recall} & \textbf{F1 Score} \\
\midrule
REAL  & 0.9835 & 0.9930 & 0.9882 \\
FAKE  & 0.9936 & 0.9849 & 0.9892 \\
\bottomrule
\end{tabular}
\end{table}

\begin{table}[H]
\centering
\caption{Confusion matrix (test set: 8{,}980 articles)}
\label{tab:confusion}
\begin{tabular}{@{}lcc@{}}
\toprule
                      & \textbf{Predicted REAL} & \textbf{Predicted FAKE} \\
\midrule
\textbf{Actual REAL}  & 4{,}240                & 30 \\
\textbf{Actual FAKE}  & 71                     & 4{,}639 \\
\bottomrule
\end{tabular}
\end{table}

The model achieves near-perfect accuracy on the ISOT dataset. However, it
is important to note that this performance reflects the dataset's
characteristics (clear stylistic differences between the real and fake
article sources) and may not generalise to adversarial or out-of-domain
content.

\begin{figure}[H]
    \centering
    \includegraphics[width=0.55\textwidth]{confusion_matrix.png}
    \caption{Confusion matrix for the Logistic Regression classifier on the
             ISOT test set.}
    \label{fig:confusion}
\end{figure}


\subsection{Conditional RAG Behaviour}

The 85\% confidence threshold was chosen based on empirical observation:
articles with ML confidence above this threshold typically receive
consistent verdicts from the agents regardless of retrieved evidence. For
articles below this threshold, the RAG evidence meaningfully influences
agent reasoning by providing ground-truth reference articles for
comparison.

\subsection{Multi-Agent Agreement Patterns}

Qualitative analysis of the system's outputs reveals the following patterns:

\begin{itemize}[leftmargin=2em]
    \item \textbf{High agreement (3/3 aligned)}: Observed on clearly fake
          or clearly real articles where ML confidence is high and RAG
          evidence is consistent.
    \item \textbf{Medium agreement (2/1 split)}: Typically occurs on
          articles with ambiguous language or where the ML confidence is
          moderate (60--85\%). Agent~B (Sceptical) is most likely to
          dissent.
    \item \textbf{Low agreement}: Rare; generally indicates genuinely
          ambiguous content where the system appropriately signals
          uncertainty.
\end{itemize}


% ════════════════════════════════════════════════════════════
\section{Discussion}
% ════════════════════════════════════════════════════════════

\subsection{Strengths}

\begin{enumerate}[leftmargin=2em]
    \item \textbf{Interpretability}: Unlike a black-box classifier, the
          system provides three independent reasoning perspectives and a
          synthesis explanation, allowing users to understand \emph{why} a
          verdict was reached.
    \item \textbf{Layered confidence}: The combination of ML confidence,
          agent agreement level, and computed risk factors gives users
          multiple axes for calibrating trust.
    \item \textbf{Efficient deployment}: Precomputed embeddings and an
          ephemeral vector store eliminate the need for persistent storage
          or GPU inference, enabling deployment on free-tier cloud services.
    \item \textbf{Conditional RAG}: Skipping retrieval for high-confidence
          predictions reduces average latency without sacrificing quality
          on ambiguous inputs.
    \item \textbf{Dual-provider LLM}: Using separate LLM providers (Groq
          for agents, NVIDIA for the judge) reduces single-provider bias
          and adds resilience through fallback.
\end{enumerate}

\subsection{Limitations}

\begin{enumerate}[leftmargin=2em]
    \item \textbf{Dataset bias}: The ISOT dataset contains articles from
          specific sources (Reuters for real, Politifact-flagged sites for
          fake). The ML model and RAG evidence may not generalise well to
          articles from other domains, languages, or writing styles.
    \item \textbf{Sequential agent execution}: Agents run in series
          (A~$\to$~B~$\to$~C), adding latency. Parallel execution would
          reduce response time but requires careful state management.
    \item \textbf{Cold start}: On Render's free tier, the first request
          after a sleep period takes 15--30 seconds due to model loading
          and embedding initialisation.
    \item \textbf{Single-mode backend}: The frontend offers a ``Classical
          ML'' mode, but the backend always runs the full agentic pipeline.
          A dedicated ML-only endpoint would reduce latency for users who
          do not need agent reasoning.
    \item \textbf{Embedding sample size}: Only 5{,}000 of the
          approximately 44{,}000 ISOT articles are embedded for RAG
          retrieval, limiting evidence coverage.
    \item \textbf{No real-time fact-checking}: The system analyses
          linguistic patterns and compares against a static dataset. It
          does not perform live web searches or consult fact-checking
          databases.
\end{enumerate}


% ════════════════════════════════════════════════════════════
\section{Responsible AI Considerations}
% ════════════════════════════════════════════════════════════

The design of this system incorporates several responsible AI principles:

\begin{enumerate}[leftmargin=2em]
    \item \textbf{Transparency}: Every verdict is accompanied by
          per-agent reasoning, evidence references, and risk factors.
          Users can inspect the full chain of reasoning, not just a final
          label.
    \item \textbf{Uncertainty signalling}: When the system is uncertain
          (low ML confidence, low agent agreement, mixed RAG evidence),
          risk factors are explicitly surfaced rather than hidden.
    \item \textbf{Anti-hallucination measures}: LLM agents are prompted
          with ``Respond ONLY based on the evidence provided. Never invent
          or hallucinate facts.'' RAG evidence is drawn from a verified
          dataset.
    \item \textbf{Multi-perspective reasoning}: By design, the three
          agents represent different analytical stances, reducing the risk
          of a single-perspective bias dominating the verdict.
    \item \textbf{Disclaimer-ready design}: The system is designed as a
          decision-support tool, not an authority. Risk factors and
          confidence levels encourage users to exercise their own judgment.
    \item \textbf{Data privacy}: Article text is processed in real time
          and not stored. No user data is persisted on the server.
    \item \textbf{API key security}: Secrets are managed via environment
          variables and are never committed to the repository. A
          \texttt{.env.example} template documents required variables
          without exposing real values.
\end{enumerate}


% ════════════════════════════════════════════════════════════
\section{Conclusion}
% ════════════════════════════════════════════════════════════

This project demonstrates that combining classical machine learning,
retrieval-augmented generation, and multi-agent LLM reasoning into a
unified pipeline produces a news credibility analysis system that is both
accurate and interpretable. The ML layer provides a fast, high-accuracy
initial signal; the RAG layer grounds the analysis in verified evidence;
and the multi-agent system delivers transparent, multi-perspective
reasoning that enables users to understand and evaluate the verdict.

The conditional RAG mechanism and lazy-loading architecture make the system
deployable on free-tier cloud infrastructure, demonstrating that
sophisticated AI pipelines need not require expensive compute resources.

Future work could extend the system with parallel agent execution for lower
latency, integration with live fact-checking APIs for real-time evidence,
expansion of the embedding corpus beyond the ISOT dataset, and a dedicated
ML-only endpoint for latency-sensitive use cases.


% ════════════════════════════════════════════════════════════
% Appendices (optional)
% ════════════════════════════════════════════════════════════

\newpage
\appendix

\section{Frontend Screenshots}
\label{app:screenshots}

\begin{figure}[H]
    \centering
    \includegraphics[width=0.9\textwidth]{landing_page.png}
    \caption{Article input interface with mode selection.}
    \label{fig:screenshot-input}
\end{figure}

\begin{figure}[H]
    \centering
    \includegraphics[width=0.9\textwidth]{thinking.png}
    \caption{Real-time processing pipeline showing active analysis steps.}
    \label{fig:screenshot-pipeline}
\end{figure}

\begin{figure}[H]
    \centering
    \includegraphics[width=0.9\textwidth]{result.png}
    \caption{Final verdict display with agent reasoning, evidence summary,
             and risk indicators.}
    \label{fig:screenshot-verdict}
\end{figure}

\end{document}
